\section{Relation to Existing Theory and Research Program}
\label{sec:discussion}
The framework keeps Shannon's uncertainty accounting and adds a response-resolved organization of prediction. The operator \(\G_c\) belongs to the conditional-estimation and principal-inertia tradition~\citep{calmon2017}. Restricted prediction classes give the computational perspective developed by \citet{xu2020}. Equivalence classes of histories with common future implications connect the continuation view to computational mechanics~\citep{shalizi2001}. Our emphasis is the combination of these predictive distinctions with a registration law that determines their phase under composition.

Geometric phase supplies the appropriate comparison mathematics. Local alignment can accumulate into a finite relative orientation when several contexts are joined~\citep{pancharatnam1956,berry1984,simon1983}. In the prediction setting this orientation identifies which future question a transferred response addresses. \Cref{eq:phaseloss} gives a direct statistical consequence. The matrix-valued formulation handles several modes, while a two-mode sector supplies the familiar complex plane.

The informationon hypothesis combines spinor interval geometry~\citep{penrose1960} and causal ordering~\citep{surya2019}. Its two-sided amplitude law separates coordinate changes from decompositions with a shared completion. Physical tests require a dynamical specification of the covariance, event-graph weights and couplings, followed by comparison with measured propagation and correlation statistics.

The next predictive test is equally concrete. Calibrate a family of future contrasts and response routes, estimate their mode strengths and phases, and predict the consequences of changing a route or its registration. During learning, track how those quantities change as additional evidence becomes usable. Stable spectral bands permit compact estimates; changes in rank or weak overlap identify where a new mode description is required. The error bounds in \cref{app:errors} specify the accuracy needed for the resulting predictions.

\section{Conclusion}
Synchronicity is the common organization through which present records constrain future outcomes. Conditional prediction measures its strength; information phase records the relative orientation of predictive distinctions carried through different response routes. The rhombic geometry joins parallel participation and ordered recombination. Learning reuses these relations, and the informationon construction proposes their physical realization.
